\begin{frame}{Multi-hop}
\framesubtitle{Nodes help each other to reach the base station}
\begin{columns}
\begin{column}{0.6\textwidth}
\begin{itemize}
    % \item Communication between neighbor
    \item Routing protocol needed to reach destination
    \begin{itemize}
      \item Range increased
      \item Adapts topology to node failure
    \end{itemize}
    \item Increase reliability
    % \item Bi-directional
    \item Overall higher energy consumption 
\end{itemize}
\end{column}
\begin{column}{0.4\textwidth}
\begin{figure}[H]
    \centering
    \scalebox{0.8}{
    \begin{tikzpicture}[auto, thick]
      % Place super peers and connect them
      \foreach \place/\name in {{(0,0)/a}}
        \node[gateways] (\name) at \place {};
      \node[motes] (a4) at (-1.7, 1) {};
      \node[motes] (a5) at (-2.8, 1.3) {};
      \node[motes] (a7) at (-2.5, -0.5) {};
      \node[motes] (a8) at (-1.8, -1) {};
     \foreach \pos/\i in {below left of/1, below of/2, left of/3}
        \node[motes, \pos =a ] (a\i) {};
      \foreach \speer/\peer in {a/a1,a/a2,a/a3,a7/a3,a8/a3}
        \path[dotted] (\speer) edge (\peer);

      \node[motes, above left of=a ] (a6) {};
      \path[dotted,->] (a5) edge (a4);
      \path[dotted,->] (a4) edge (a6);
      \path[dotted,->] (a6) edge (a);
      \draw[dotted,draw={black},color=gray!70] (-2.8,1.3) circle (1.3cm);
      \draw[dotted,draw={black},color=gray!70] (-1.65,1) circle (1.3cm);
      \draw[dotted,draw={black},color=gray!70] (-0.7,0.7) circle (1.3cm);
    \end{tikzpicture}
    }
    \caption{Multi-hop Communication}
    \scalebox{0.5}{%
    \begin{tabular}{r@{: }l r@{: }l}
    \begin{tikzpicture}\draw[left color=orange,line width=1pt] circle(1ex);\end{tikzpicture} & Gateway & \begin{tikzpicture}\draw[left color=gray,line width=1pt] circle(1ex);\end{tikzpicture} & Motes
    \end{tabular}
    }
\end{figure}
\end{column}
\end{columns}
\end{frame}

% \begin{frame}{Multi-hop}
% \framesubtitle{LoRaWAN Extension}
% \begin{figure}[H]
%     \centering
%     \includegraphics[width=0.98\textwidth]{presentation.tex/fig/lorawanextension.png}
%     \caption{LoRaWAN Extension\footnotemark}
% \end{figure}

% \footcitetext{DIAS2018424}
% \end{frame}

% \begin{frame}{Multi-hop}
% \framesubtitle{Linear Sensor Network}
% \begin{columns}
% \begin{column}{0.5\textwidth}
% \begin{figure}[H]
%     \centering
%     \includegraphics[width=1\textwidth]{presentation.tex/fig/lsnlora.jpg}
%     \caption{Underground tunnel installation\footnotemark}
% \end{figure}
% \end{column}
% \begin{column}{0.5\textwidth}
% \begin{figure}[H]
%     \centering
%     \includegraphics[width=0.98\textwidth]{presentation.tex/fig/lsnlora2.jpg}
%     \caption{Protocol\footnotemark}
% \end{figure}
% \end{column}
% \end{columns}
% \footcitetext{Abrardo_2019}
% \end{frame}



% \begin{frame}{Multi-hop}
% \framesubtitle{LoRaBlink}
% \begin{columns}
% \begin{column}{0.3\textwidth}
% \begin{figure}[H]
%     \centering
%     \includegraphics[width=1\textwidth]{presentation.tex/fig/lorablink2.png}
%     \caption{Range\footnotemark}
% \end{figure}
% \end{column}
% \begin{column}{0.7\textwidth}
% \begin{figure}[H]
%     \centering
%     \includegraphics[width=0.98\textwidth]{presentation.tex/fig/lorablink.png}
%     \caption{LoRaBlink protocol\footnotemark}
% \end{figure}
% \end{column}
% \end{columns}
% \footcitetext{lorablink}
% \end{frame}

% \begin{frame}{Multi-hop}
% \framesubtitle{LoRa Multi-hop Network}
% % Sum up of multi-hop LoRa
% % What related work achieved
% \begin{itemize}
%     \item Cope with collision through
%     \begin{itemize}
%         \item Node coordination
%         \item Clock synchronization
%     \end{itemize}
%     \item Cope with bordering nodes through
%     \begin{itemize}
%         \item Increased range
%     \end{itemize}
%     \item Linear Sensor Network
%     % \item Interference resilience
% \end{itemize}
% \end{frame}

% \begin{frame}{Multi-hop}
% \framesubtitle{LoRa Multi-hop Network}
% % My work what it will achieve based on that
% \begin{itemize}
%     \item Low-power
%     \begin{itemize}
%         \item Time-Division
%         \item Synchronization
%         \item Scheduling
%     \end{itemize}
%     \item Reliable
%     \begin{itemize}
%         \item No collision
%         \item Interference resistance
%         \item Frequency-Division
%     \end{itemize}
%     \item Bi-directional
% \end{itemize}
% \end{frame}

% \begin{frame}{Multi-hop}
% \framesubtitle{Solving problem 1}
% \begin{figure}[H]
%     \centering
%     \scalebox{1}{%
%   \begin{tikzpicture}[]

%   \begin{scope}[
%       xshift=-0.2cm,
%       asn/.style={black!70, minimum width=2cm},
%       timeslot/.style={draw, rectangle, minimum width=2cm, minimum height=1cm},
%       arr/.style={help lines,black!70,<->},
%   ]
%     \foreach \i in {0,...,4} {
%       \node (ts\i) [asn] at (2*\i, 1) {$TS_{\i}$};
%     }
%     \node (tss0) [timeslot] at (0, 0) {\tiny Shared};
%     \node (tss1) [timeslot] at (2, 0) {\tiny B $\rightarrow$ A};
%     \node (tss2) [timeslot] at (4, 0) {\tiny C $\rightarrow$ B};
%     \node (tss3) [timeslot] at (6, 0) {\tiny C $\rightarrow$ A};
%     \node (tss4) [timeslot] at (8, 0) {\tiny A $\rightarrow$ C};
%     \path[->] (-1.2,-1) edge (9.5,-1);
%     \node [] at (10, -1.2) {time};
%   \end{scope}
%   \begin{scope}[xshift=3.8cm,yshift=-2cm,->,>=stealth',shorten >=1pt,auto,node distance=2.5cm]
%     \tikzstyle{every state}=[thick,draw=gray!50,fill=gray!20,draw=none,text=black]

%     \node[state]         (A) [] {A};
%     \node[state]         (B) [below of=A]       {B};
%     \node[state]         (C) [right of=B]       {C};

%     \path (A) edge [bend left] node {$TS_4$} (C)
%           (B) edge [bend left] node {$TS_1$} (A)
%           (C) edge [bend left] node[above right] {$TS_3$} (A)
%               edge [bend left] node {$TS_2$} (B);
%   \end{scope}
% \end{tikzpicture}
% }
% \end{figure}


% \end{frame}
